% Maurice Heide, public CV for https://mooofna.github.io/cv/
% Derived from the application CV of 09.10.2026 with the phone number and patient counts removed
% and a link to the homepage added (it links GitHub and Google Scholar). Ranges are written with "to", not dashes.
% Build: npm run cv (xelatex via latexmk; the Libertine and Biolinum fonts ship with TeX Live).
\documentclass[10pt,a4paper]{article}
\usepackage[a4paper,top=1.0cm,bottom=1.0cm,left=1.6cm,right=1.6cm]{geometry}
\usepackage[english]{babel}
\usepackage{fontspec}
\usepackage{xcolor}
\usepackage{microtype}
\usepackage{enumitem}
\usepackage{ragged2e}
\usepackage[hidelinks]{hyperref}

% Fonts by file name (they ship with TeX Live), so the build does not depend on installed system fonts.
\setmainfont{LinLibertine}[Extension=.otf,UprightFont=*_R,BoldFont=*_RZ,ItalicFont=*_RI,BoldItalicFont=*_RZI,
  Ligatures={TeX,Common},Numbers={Lining,Proportional}]
\newfontfamily\ui{LinBiolinum}[Extension=.otf,UprightFont=*_R,BoldFont=*_RB,ItalicFont=*_RI,
  Ligatures=TeX,Numbers=Lining]
\definecolor{akzent}{HTML}{17365D}
\definecolor{grau}{HTML}{4A4A4A}
\definecolor{linie}{gray}{0.62}

\hypersetup{pdftitle={Maurice Heide, CV},pdfauthor={Maurice Heide},pdfcreator={Maurice Heide},pdfproducer={}}

\pagestyle{empty}
\setlength{\parindent}{0pt}
\setlength{\parskip}{0pt}
\hyphenpenalty=10000 \exhyphenpenalty=10000 \emergencystretch=3em
\newcommand{\grundtext}{\fontsize{10.1}{12.1}\selectfont}
\AtBeginDocument{\grundtext}

\newcommand{\trenner}{\unskip\hspace{0.7em}\textbar\hspace{0.7em}\ignorespaces}
\newcommand{\sperr}[1]{{\addfontfeature{LetterSpace=8}#1}}

% Abschnitt: Biolinum halbfett, Versalien gesperrt, Akzentfarbe, feine Linie
\newcommand{\abschnitt}[1]{%
  \par\vspace{4.6pt}%
  {\ui\bfseries\fontsize{10.5}{10.5}\selectfont\color{akzent}\sperr{\MakeUppercase{#1}}}\par
  \vspace{2.2pt}{\color{linie}\hrule height 0.4pt}\vspace{4pt}}

% Eintrag: Institution fett links, Ort rechts; Rolle kursiv links, Zeitraum rechts
\newcommand{\eintrag}[4]{%
  \par\vspace{2.4pt}%
  \noindent\textbf{#1}\hfill{\ui\fontsize{9.5}{12.6}\selectfont\color{grau}#2}\par
  \noindent\textit{#3}\hfill{\ui\fontsize{9.5}{12.6}\selectfont\color{grau}#4}\par}

% Stichpunkte
\setlist[itemize]{label={\fontsize{7}{7}\selectfont\textbullet},leftmargin=1.05em,labelsep=0.45em,
  topsep=0.8pt,itemsep=0.6pt,parsep=0pt,partopsep=0pt,before=\RaggedRight}

% Zeile im Skills- und Awards-Block: Etikett fett, Text daneben
\newcommand{\zeile}[2]{\par\noindent\hangindent=7.1em\hangafter=1
  \makebox[7.1em][l]{\textbf{#1}}#2\par\vspace{1.2pt}}

\begin{document}

% ------------------------------------------------------------------ Kopf
\begin{center}
  {\ui\bfseries\fontsize{21}{22}\selectfont\color{akzent}\sperr{Maurice Heide}}\par
  \vspace{4pt}
  {\ui\fontsize{9.8}{12}\selectfont\color{grau}%
  Zurich, Switzerland\trenner
  \href{mailto:mheide@student.ethz.ch}{mheide@student.ethz.ch}\trenner
  \href{https://mooofna.github.io}{mooofna.github.io}\trenner
  \href{https://www.linkedin.com/in/maurice-heide}{linkedin.com/in/maurice-heide}}\par
\end{center}
\vspace{-3pt}

% ------------------------------------------------------------------ Ausbildung
\abschnitt{Education}

\eintrag{ETH Zurich}{Zurich, Switzerland}{M.Sc.\ Data Science}{Sep 2026 to present}
\begin{itemize}
\item Coursework: Probabilistic AI, Mathematics of Data Science, Big Data, Seminar in Representation Learning
\end{itemize}

\eintrag{TU Dresden}{Dresden, Germany}{B.Sc.\ Computer Science, grade 1.0 (top grade), passed with distinction}{Oct 2022 to Aug 2026}
\begin{itemize}
\item Thesis (1.0) on radiology foundation models at Harvard Medical School and Massachusetts General Hospital
\item All four mathematics exams 1.0 (30 ECTS); teaching assistant for Discrete Mathematics, best-rated group
\end{itemize}

\eintrag{Nanyang Technological University}{Singapore}{Exchange semester, GPA 4.83/5.00; the only exchange student of the relaunched partnership}{Jan 2025 to May 2025}

\eintrag{Seoul National University}{Seoul, South Korea}{Exchange semester, GPA 4.2/4.3}{Sep 2024 to Dec 2024}
\begin{itemize}
\item One of two exchange places; A+ in Deep Learning (graduate level) and Statistics
\end{itemize}

% ------------------------------------------------------------------ Erfahrung
\abschnitt{Experience}

\eintrag{Harvard Medical School \& Massachusetts General Hospital}{Boston, USA}{Visiting Undergraduate Researcher, Harvard University; Fintelmann Lab, Radiology}{May 2026 to Sep 2026}
\begin{itemize}
\item Bachelor thesis: benchmarked six published CT foundation models on 11 cohorts (20,862 scans) and distilled the two strongest into 14 compact 3D models (0.04M to 44M parameters)
\item Built the training and evaluation stack: PyTorch DDP on up to 8 A100s with Slurm auto-resume (2,260 GPU-hours) and a tested evaluation package whose outputs match the original models byte for byte
\item A 3M student is within 0.2 AUROC points of its 79M teacher (MLP probe), and the smallest runs 82.6$\times$ faster on CPU
\item As the lab's only engineer, scaled its CT biomarker pipeline to 3,864 scans; checked its outputs against the images, caught and removed wrong values in 1,149 scans, and set the quality-control rules
\item Took an LLM pipeline from design to production (Qwen3-4B on vLLM); it labels lung-cancer recurrence in clinical notes, with a review GUI for the medical team
\end{itemize}

\eintrag{Else Kröner Fresenius Center for Digital Health, TU Dresden}{Dresden, Germany}{Research Assistant, Kather Lab}{Oct 2025 to Mar 2026, Nov 2023 to Jun 2024}
\begin{itemize}
\item Shared first author of MuCoDi: distilled three pathology foundation models of up to 1.1B parameters into edge models; a 6.4M student, 99$\times$ smaller than the best teacher, is within 0.9 AUROC points of it on external data, and sub-million students run up to 605$\times$ faster on a Raspberry Pi 5
\item Built a WebAssembly engine that predicts biomarkers from gigapixel slides in seconds, privately in the browser
\item 2023 to 2024: turned 16\,TB of whole-slide images into 23\,MB of embeddings with a self-supervised pipeline
\end{itemize}

\eintrag{EY \& Associés, Forensic \& Integrity Services}{Paris, France}{Data Science Intern}{Jul 2025 to Sep 2025}
\begin{itemize}
\item Built a detector for AI-generated forgeries in EY's document-analysis platform (500+ clients, 25+ countries): compared 20 pretrained models, then fine-tuned DINOv3, using SVD to separate image content from forgery traces
\item Curated a balanced training set of 250k real and 250k synthetic images from 6M, controlling for the JPEG and image-size biases that weaken detectors on unseen generators
\item 99.8\% accuracy on a public benchmark of 25 generators, 97 to 100\% on passports and receipts forged by Nano Banana
\end{itemize}

\eintrag{Fraunhofer IVI}{Dresden, Germany}{Research Assistant, Applied Machine Learning}{Mar 2023 to Sep 2023}
\begin{itemize}
\item Built an automated calibration framework for the 29 near-infrared classifiers of a waste-sorting plant; XGBoost on 58 spectral parameters raised sorting accuracy by 22\% for Germany's leading waste-management company
\item Wrote an ETL pipeline whose statistical thresholds flag sensor drift and calibration errors, reducing plant downtime
\end{itemize}

% ------------------------------------------------------------------ Publikation
\abschnitt{Publication}

\noindent T.~Lenz\textsuperscript{\textdagger}, \textbf{M.~Heide}\textsuperscript{\textdagger}, M.~Gustav, N.~G.~Reitsam, J.~N.~Kather. \textit{Multi-Teacher Contrastive Distillation for Edge-Efficient Pathology Foundation Models.} MICCAI 2026 Workshop on Efficient Medical AI. \href{https://arxiv.org/abs/2607.05533}{arXiv:2607.05533}. \textsuperscript{\textdagger}Equal contribution.\par

% ------------------------------------------------------------------ Auszeichnungen
\abschnitt{Honors}

\begin{itemize}
\item Best Paper Award, MICCAI 2026 Workshop on Efficient Medical Artificial Intelligence (EMA4MICCAI), Strasbourg
\item Scholar of the German Academic Scholarship Foundation (Studienstiftung des deutschen Volkes) since 2022
\item Hermann Medal for the best Abitur of the year (1.0); two prizes of the German Physical Society
\end{itemize}

% ------------------------------------------------------------------ Kenntnisse
\abschnitt{Skills}

\zeile{Programming}{Python (PyTorch, JAX, scikit-learn, XGBoost, pandas, Hugging Face, ONNX), SQL, Rust, C++, Java}
\zeile{Open source}{\href{https://github.com/mooofna/turing-machines}{turing-machines} (Rust, WebAssembly): decides all 858,909 four-state Turing machines}
\zeile{Infrastructure}{Slurm and Apptainer on GPU clusters, Docker, Kubernetes, GitHub Actions, pytest, mypy, W\&B}
\zeile{Languages}{German (native), English (C2, IELTS 8.5), French (A1)}

\end{document}
